\section{Introduction}
\label{sec:intro}

LLM coding agents increasingly resolve real GitHub issues by
editing a repository directly, and benchmarks such as SWE-Bench
now track how often they succeed \citep{jimenez2024swebench}. A
patch can only be correct, however, if it changes the right
code. Before any edit, the agent must identify the files and
functions an issue actually touches, the step known as
\emph{code localization}. This step is decisive and largely
unrecoverable: once a wrong location enters the agent's editable
context, the repair stage reasons over the wrong code, and no
later reasoning recovers the correct patch
\citep{jimenez2024swebench, yang2024sweagent}. \citet{xia2024agentless}
drop agentic interaction and resolve issues with a fixed
localize-then-repair pipeline, crediting much of their accuracy
to a hierarchical localization phase, and methods that improve
the localization stage report corresponding gains in end-to-end
resolution \citep{sohrabizadeh2025nemotroncortexa}. Localization
is therefore not a preprocessing convenience but the step that
determines how well downstream repair can perform.

% Yet localizing well is hard: a real repository holds thousands
% of files joined by tangled import, call, and inheritance edges,
% while an issue report describes symptoms in natural language
% that can sit far from the code that must change. 
Current systems close this gap in four ways. Reinforcement learning methods
train the search policy itself: RepoSearcher pairs
tool-integrated RL with rejection-sampled supervised
fine-tuning \citep{ma2025tooltrain}, and SWE-RL scales RL over
open software-evolution data \citep{wei2025swerl}.
General-purpose agent scaffoldings give a capable LLM an
open-ended action space with shell access, as in SWE-Agent and
OpenHands \citep{yang2024sweagent, wang2025openhands}.
Graph-exploration agents such as LocAgent let an LLM traverse a
repository code graph, hopping freely between related code units
\citep{chen2025locagent}. Dense retrieval methods skip the agent
altogether: SWERank scores code units with a bi-encoder
retriever and a listwise LLM scorer
\citep{reddy2025sweranksoftwareissuelocalization}.

The first three families pay for accuracy in tokens, and none of
them bounds the cost. An RL agent, an open-ended scaffold, and a
graph explorer all keep issuing tool calls and appending
observations until they choose to stop, so a localization
trajectory grows without a predictable ceiling. Under an
OpenHands-style scaffolding, a single localization run can
exceed 131K tokens. A model with a
long context window and ample compute absorbs that cost; the
open 7--14B code models most teams can realistically run do not.
A Qwen-class coder with a 32K-token window overflows before the
search terminates. Reinforcement learning compounds the problem:
a multi-step episode returns only a sparse terminal reward, so
adapting a small model to it takes many long and expensive
rollouts. Dense retrieval is the one cheap option, but a
single-step, order-centric ranker cannot revisit a discarded
candidate or follow a structural edge, so it leaves the
sequential, graph-structured side of the problem untouched. No
existing method is at once bounded in token cost at inference
and cheap enough to train on a small open model.

% ============================================================
% OUTLINE -- Introduction paragraph 4 (CONTRIBUTIONS) -- TO WRITE
% Close on explicit contributions + a one-line results preview,
% then transition to Sec. 2. Draft pending user decision on the
% contribution list.
%   C1  a multi-round, bounded-context localization scaffolding
%       for small code LLMs;
%   C2  a telescoping result: the dense per-round target shares
%       its optimum with end-to-end recall@L (no myopia gap),
%       converting a sparse-reward multi-step RL problem into
%       supervised learning;
%   C3  a supervised objective -- weighted BCE + a differentiable
%       coverage term -- provably co-optimal with recall@L, and
%       the set-coverage reframing it rests on;
%   C4  results: small open models competitive with GPT Codex
%       5.3 under the same scaffolding; Qwen-30B competitive
%       with OpenHands at a fraction of the token cost; on
%       SWE-Bench Lite, SWE-Bench Verified, and SWE-PolyBench.
% ============================================================
